\section{AMC 10 10A 2015}
        \begin{enumerate}
        	\item (\textit{AMC 10 10A 2015}) What is the value of $(2^0-1+5^2+0)^{-1}\times5?$


 $\textbf{(A)}\ -125\qquad\textbf{(B)}\ -120\qquad\textbf{(C)}\ \frac{1}{5}\qquad\textbf{(D)}\ \frac{5}{24}\qquad\textbf{(E)}\ 25$


 

	\item (\textit{AMC 10 10A 2015}) A box contains a collection of triangular and square tiles. There are $25$ tiles in the box, containing $84$ edges total. How many square tiles are there in the box?


 $\textbf{(A)}\ 3\qquad\textbf{(B)}\ 5\qquad\textbf{(C)}\ 7\qquad\textbf{(D)}\ 9\qquad\textbf{(E)}\ 11$


 

	\item (\textit{AMC 10 10A 2015}) Ann made a 3-step staircase using 18 toothpicks as shown in the figure.  How many toothpicks does she need to add to complete a 5-step staircase?


 
                    \begin{figure}[H]
                        \centering
                        \includegraphics[width=0.5\textwidth]{Problem/AMC_10_10A_2015_3_Problem_Figure_1.png}
                    \end{figure}
                    


 $\textbf{(A)}\ 9\qquad\textbf{(B)}\ 18\qquad\textbf{(C)}\ 20\qquad\textbf{(D)}\ 22\qquad\textbf{(E)}\ 24$


 

	\item (\textit{AMC 10 10A 2015}) Pablo, Sofia, and Mia got some candy eggs at a party.  Pablo had three times as many eggs as Sofia, and Sofia had twice as many eggs as Mia.  Pablo decides to give some of his eggs to Sofia and Mia so that all three will have the same number of eggs.  What fraction of his eggs should Pablo give to Sofia?


 $\textbf{(A)}\ \frac{1}{12}\qquad\textbf{(B)}\ \frac{1}{6}\qquad\textbf{(C)}\ \frac{1}{4}\qquad\textbf{(D)}\ \frac{1}{3}\qquad\textbf{(E)}\ \frac{1}{2}$


 

	\item (\textit{AMC 10 10A 2015}) Mr. Patrick teaches math to $15$ students. He was grading tests and found that when he graded everyone's test except Payton's, the average grade for the class was $80$. After he graded Payton's test, the test average became $81$. What was Payton's score on the test?


 $\textbf{(A)}\ 81\qquad\textbf{(B)}\ 85\qquad\textbf{(C)}\ 91\qquad\textbf{(D)}\ 94\qquad\textbf{(E)}\ 95$


 

	\item (\textit{AMC 10 10A 2015}) The sum of two positive numbers is $5$ times their difference. What is the ratio of the larger number to the smaller number?


 $\textbf{(A)}\ \frac{5}{4}\qquad\textbf{(B)}\ \frac{3}{2}\qquad\textbf{(C)}\ \frac{9}{5}\qquad\textbf{(D)}\ 2 \qquad\textbf{(E)}\ \frac{5}{2}$


 

	\item (\textit{AMC 10 10A 2015}) How many terms are there in the arithmetic sequence $13$, $16$, $19$, . . ., $70$, $73$?


 $\textbf{(A)}\ 20\qquad\textbf{(B)}\ 21\qquad\textbf{(C)}\ 24\qquad\textbf{(D)}\ 60\qquad\textbf{(E)}\ 61$


 

	\item (\textit{AMC 10 10A 2015}) Two years ago Pete was three times as old as his cousin Claire. 2 years before that, Pete was four times as old as Claire. In how many years will the ratio of their ages be $2$ : $1$?


 $\textbf{(A)}\ 2\qquad\textbf{(B)}\ 4\qquad\textbf{(C)}\ 5\qquad\textbf{(D)}\ 6\qquad\textbf{(E)}\ 8$


 

	\item (\textit{AMC 10 10A 2015}) Two right circular cylinders have the same volume. The radius of the second cylinder is $10\%$ more than the radius of the first. What is the relationship between the heights of the two cylinders?


 $\textbf{(A)}\  \text{The second height is } 10\% \text{ less than the first.}$
$\textbf{(B)}\  \text{The first height is } 10\% \text{ more than the second.}$
$\textbf{(C)}\ \text{The second height is } 21\% \text{ less than the first.}$
$\textbf{(D)}\ \text{The first height is } 21\% \text{ more than the second.}$
$\textbf{(E)}\ \text{The second height is } 80\% \text{ of the first.}$


 

	\item (\textit{AMC 10 10A 2015}) How many rearrangements of $abcd$ are there in which no two adjacent letters are also adjacent letters in the alphabet?  For example, no such rearrangements could include either $ab$ or $ba$.


 $\textbf{(A)}\ 0\qquad\textbf{(B)}\ 1\qquad\textbf{(C)}\ 2\qquad\textbf{(D)}\ 3\qquad\textbf{(E)}\ 4$


 

	\item (\textit{AMC 10 10A 2015}) The ratio of the length to the width of a rectangle is $4$ : $3$. If the rectangle has diagonal of length $d$, then the area may be expressed as $kd^2$ for some constant $k$. What is $k$?


 $\textbf{(A)}\ \frac{2}{7}\qquad\textbf{(B)}\ \frac{3}{7}\qquad\textbf{(C)}\ \frac{12}{25}\qquad\textbf{(D)}\ \frac{16}{25}\qquad\textbf{(E)}\ \frac{3}{4}$


 

	\item (\textit{AMC 10 10A 2015}) Points $(\sqrt{\pi}, a)$ and $(\sqrt{\pi}, b)$ are distinct points on the graph of $y^2+x^4=2x^2y+1$. What is $|a-b|$?


 $\textbf{(A) }1\qquad\textbf{(B) }\dfrac{\pi}{2}\qquad\textbf{(C) }2\qquad\textbf{(D) }\sqrt{1+\pi}\qquad\textbf{(E) }1+\sqrt{\pi}$


 

	\item (\textit{AMC 10 10A 2015}) Claudia has 12 coins, each of which is a 5-cent coin or a 10-cent coin. There are exactly 17 different values that can be obtained as combinations of one or more of her coins. How many 10-cent coins does Claudia have?


 $\textbf{(A) }3\qquad\textbf{(B) }4\qquad\textbf{(C) }5\qquad\textbf{(D) }6\qquad\textbf{(E) }7$


 

	\item (\textit{AMC 10 10A 2015}) The diagram below shows the circular face of a clock with radius $20$ cm and a circular disk with radius $10$ cm externally tangent to the clock face at $12$ o'clock. The disk has an arrow painted on it, initially pointing in the upward vertical direction. Let the disk roll clockwise around the clock face. At what point on the clock face will the disk be tangent when the arrow is next pointing in the upward vertical direction?


 
                    \begin{figure}[H]
                        \centering
                        \includegraphics[width=0.5\textwidth]{Problem/AMC_10_10A_2015_14_Problem_Figure_1.png}
                    \end{figure}
                    


 $\textbf{(A) }\mathrm{2 o'clock} \qquad\textbf{(B) }\mathrm{3 o'clock} \qquad\textbf{(C) }\mathrm{4 o'clock} \qquad\textbf{(D) }\mathrm{6 o'clock} \qquad\textbf{(E) }\mathrm{8 o'clock}$


 

	\item (\textit{AMC 10 10A 2015}) Consider the set of all fractions $\frac{x}{y},$ where $x$ and $y$ are relatively prime positive integers. How many of these fractions have the property that if both numerator and denominator are increased by $1$, the value of the fraction is increased by $10\%$?


 $\textbf{(A) }0\qquad\textbf{(B) }1\qquad\textbf{(C) }2\qquad\textbf{(D) }3\qquad\textbf{(E) }\text{infinitely many}$


 

	\item (\textit{AMC 10 10A 2015}) If $y+4 = (x-2)^2, x+4 = (y-2)^2$, and $x \neq y$, what is the value of $x^2+y^2$?


 $\textbf{(A) }10\qquad\textbf{(B) }15\qquad\textbf{(C) }20\qquad\textbf{(D) }25\qquad\textbf{(E) }\text{30}$


 

	\item (\textit{AMC 10 10A 2015}) A line that passes through the origin intersects both the line $x=1$ and the line $y=1+\frac{\sqrt{3}}{3}x$. The three lines create an equilateral triangle. What is the perimeter of the triangle?


 $\textbf{(A) }2\sqrt{6}\qquad\textbf{(B) }2+2\sqrt{3}\qquad\textbf{(C) }6\qquad\textbf{(D) }3+2\sqrt{3}\qquad\textbf{(E) }6+\frac{\sqrt{3}}{3}$


 

	\item (\textit{AMC 10 10A 2015}) Hexadecimal (base-16) numbers are written using numeric digits $0$ through $9$ as well as the letters $A$ through $F$ to represent $10$ through $15$. Among the first $1000$ positive integers, there are $n$ whose hexadecimal representation contains only numeric digits. What is the sum of the digits of $n$?


 $\textbf{(A) }17\qquad\textbf{(B) }18\qquad\textbf{(C) }19\qquad\textbf{(D) }20\qquad\textbf{(E) }21$


 

	\item (\textit{AMC 10 10A 2015}) The isosceles right triangle $ABC$ has right angle at $C$ and area $12.5$. The rays trisecting $\angle ACB$ intersect $AB$ at $D$ and $E$. What is the area of $\bigtriangleup CDE$?


 $\textbf{(A) }\dfrac{5\sqrt{2}}{3}\qquad\textbf{(B) }\dfrac{50\sqrt{3}-75}{4}\qquad\textbf{(C) }\dfrac{15\sqrt{3}}{8}\qquad\textbf{(D) }\dfrac{50-25\sqrt{3}}{2}\qquad\textbf{(E) }\dfrac{25}{6}$


 

	\item (\textit{AMC 10 10A 2015}) A rectangle with positive integer side lengths in $\text{cm}$ has area $A$ $\text{cm}^2$ and perimeter $P$ $\text{cm}$. Which of the following numbers cannot equal $A+P$?


 $\textbf{(A) }100\qquad\textbf{(B) }102\qquad\textbf{(C) }104\qquad\textbf{(D) }106\qquad\textbf{(E) }108$


 \textbf{Note:} As it originally appeared in the AMC 10, this problem was stated incorrectly and had no answer; it has been modified here to be solvable. This is the original question:


 A rectangle with side lengths in $\text{cm}$ has an area of integer $A$ $\text{cm}^2$ and a perimeter of integer $P$ $\text{cm}$. Which of the following numbers cannot equal $A+P$?


 $\textbf{(A) }100\qquad\textbf{(B) }102\qquad\textbf{(C) }104\qquad\textbf{(D) }106\qquad\textbf{(E) }108$


 

	\item (\textit{AMC 10 10A 2015}) Tetrahedron $ABCD$ has $AB=5$, $AC=3$, $BC=4$, $BD=4$, $AD=3$, and $CD=\frac{12}5\sqrt2$.  What is the volume of the tetrahedron?


 $\textbf{(A) }3\sqrt2\qquad\textbf{(B) }2\sqrt5\qquad\textbf{(C) }\dfrac{24}5\qquad\textbf{(D) }3\sqrt3\qquad\textbf{(E) }\dfrac{24}5\sqrt2$


 

	\item (\textit{AMC 10 10A 2015}) Eight people are sitting around a circular table, each holding a fair coin. All eight people flip their coins and those who flip heads stand while those who flip tails remain seated. What is the probability that no two adjacent people will stand?


 $\textbf{(A)}\dfrac{47}{256}\qquad\textbf{(B)}\dfrac{3}{16}\qquad\textbf{(C) }\dfrac{49}{256}\qquad\textbf{(D) }\dfrac{25}{128}\qquad\textbf{(E) }\dfrac{51}{256}$


 

	\item (\textit{AMC 10 10A 2015}) The zeroes of the function $f(x)=x^2-ax+2a$ are integers. What is the sum of the possible values of $a$?


 $\textbf{(A) }7\qquad\textbf{(B) }8\qquad\textbf{(C) }16\qquad\textbf{(D) }17\qquad\textbf{(E) }18$


 

	\item (\textit{AMC 10 10A 2015}) For some positive integers $p$, there is a quadrilateral $ABCD$ with positive integer side lengths, perimeter $p$, right angles at $B$ and $C$, $AB=2$, and $CD=AD$.  How many different values of $p<2015$ are possible?


 $\textbf{(A) }30\qquad\textbf{(B) }31\qquad\textbf{(C) }61\qquad\textbf{(D) }62\qquad\textbf{(E) }63$


 

	\item (\textit{AMC 10 10A 2015}) Let $S$ be a square of side length $1$.  Two points are chosen at random on the sides of $S$.  The probability that the straight-line distance between the points is at least $\frac12$ is $\frac{a-b\pi}c$, where $a$, $b$, and $c$ are positive integers with $\gcd(a,b,c)=1$.  What is $a+b+c$?


 $\textbf{(A) }59\qquad\textbf{(B) }60\qquad\textbf{(C) }61\qquad\textbf{(D) }62\qquad\textbf{(E) }63$


 

\end{enumerate}